\section{The split-cube quadratic complement in the actual common window}
\label{sec:mixed-cubic-split-cube-window}

\paragraph{The preceding original-coefficient remainder.}
Retain every datum, page, prescribed pair, field, original coefficient,
lattice, normalization and full local condition of AC-G214-window.
In particular $h\geq2$ and $T\geq2$ are integers, $N=T-1$, and
\[
 \begin{gathered}
 H_u=\mathbb Q(\zeta_{2^{h+2}},\zeta_7,\sqrt{-q},\sqrt{-d},
                   \omega,\sqrt[3]{269}),\qquad K=\mathbb Q(\omega),\\
 u^{2+8\varepsilon}\leq2^h\leq2u^{2+8\varepsilon},\qquad
 T=\lfloor u^{399}\rfloor,\qquad0<\varepsilon<1/256.
 \end{gathered}\tag{MCNW.1}
\]
Keep exactly $Z_u,J_{u,0},J_{u,\rm mix},V_u,Y_u$,
$C_5^{\rm all}(u),U_5^{\rm all}(u),C_7^{\rm all}(u),U_7^{\rm all}(u)$
and $M_{8,7}^{++}(u)=M_8(u)\sqcup M_7(u)\sqcup Q_7(u)\sqcup Q_{7,7}(u)$
of (M77W.2)--(M77W.6). In particular the entire four-per-weight
residual-order-three family at coefficient seven is retained.

Put $\Delta_8=\{\chi_{-4},\chi_8,\chi_{-8}\}$ and define
\[
 \begin{split}
 \mathscr C=\{\ell\text{ prime}:{}&
        \ell>3,\quad\ell\ne7,269,\quad\ell\equiv1\pmod3,\\
       &269^{(\ell-1)/3}\equiv1\pmod\ell,\quad
                         \ell\not\equiv1\pmod{24}\},\\
 J_C(T)&=\sum_{\substack{\ell\in\mathscr C\\\ell\leq T}}
                       \left\lfloor\frac{T-1}{\ell-1}\right\rfloor,\\
 C(u)=\{(\ell,\chi^j\delta,e):{}&
       \ell\in\mathscr C,\quad j\in\{1,2\},\quad\delta\in\Delta_8,
       \quad\delta(\ell)=-1,\\
       &1\leq e<T,\quad\ell-1\mid e\},\\
 M_C(u)&=M_{8,7}^{++}(u)\cup C(u).
 \end{split}\tag{MCNW.2}
\]
Every rational conductor-seven and auxiliary $q,d$ character factor
of a new label is trivial. The fixed quadratic characters are already
present in the original two-power factor; the auxiliary source
parameter $d$ is not a new choice of quadratic proof field here.
Define the nonnegative remainder
\[
 \begin{split}
 \mathcal R_{399}^{\rm match,cub8C,7qr,5all}(u)={}&
  \sum_{\substack{(\ell,\eta,e)\in\\ J_{u,\rm mix}\setminus M_C(u)}}
                  \frac{\log\#\mathcal S_{\ell,\eta}(e)}{b_\ell}\\
 &+\sum_{\substack{(\ell,\eta,e)\in J_{u,0}\setminus\\
       (V_u\cup Y_u\cup C_7^{\rm all}(u)\cup C_5^{\rm all}(u))}}
       \frac{\log\#\mathcal S^{\mathbb Q}_{\ell,
                             \psi_{\eta,e}^{\rm nc}}(e)}{b_\ell}.
 \end{split}\tag{MCNW.3}
\]
The whole rational noncritical sum is the one in (M77W.3).

\begin{theorem}[AC-G216-window: the complete split-cube complement]
\label{AC-G216-window}
Every triple of $C(u)$ has zero full original $K$-strict group.
It is a new retained mixed triple, and
\[
 \begin{gathered}
 C(u)\subset J_{u,\rm mix}\setminus M_{8,7}^{++}(u),\\
 C(u)\cap\bigl(Z_u\cup V_u\cup Y_u\cup C_7^{\rm all}(u)
          \cup C_5^{\rm all}(u)\cup M_{8,7}^{++}(u)\bigr)
                                                     =\varnothing,\\
 M_C(u)=M_{8,7}^{++}(u)\sqcup C(u),\\
 \mathcal R_{399}^{\rm match,cub8C,7qr,5all}(u)
                   =\mathcal R_{399}^{\rm match,cub8,7qr,5all}(u).
 \end{gathered}\tag{MCNW.4}
\]
The exact new and cumulative mixed counts are
\[
 \begin{gathered}
 |C(u)|=4J_C(T),\\
 |M_C(u)|=8J_{\rm cub}(T)+4J_7(T)+6\lfloor N/6\rfloor+4N+4J_C(T).
 \end{gathered}\tag{MCNW.5}
\]
The retained indexing set is properly reduced for $T\geq37$.
The full disjoint represented-triple omission and its size are
\[
 \begin{gathered}
 Z_u\sqcup V_u\sqcup Y_u\sqcup U_7^{\rm all}(u)
                    \sqcup U_5^{\rm all}(u)\sqcup M_C(u),\\
 |Z_u|+|V_u|+|Y_u|+51N-\lfloor N/4\rfloor
       +8J_{\rm cub}(T)+4J_7(T)+6\lfloor N/6\rfloor+4J_C(T).
 \end{gathered}\tag{MCNW.6}
\]
With exactly the preceding constants and costs,
\[
 \begin{aligned}
 |I_{399}-\mathcal R_{399}^{\rm match,cub8C,7qr,5all}|&\leq
        \mathcal E_u+\mathcal L_u+\mathcal A_u+C_2\mathcal P_u,\\
 |I_{399}^{\rm all7,2}-\mathcal R_{399}^{\rm match,cub8C,7qr,5all}|&\leq
        \mathcal E_u+\mathcal L_u+C_2\mathcal P_u.
 \end{aligned}\tag{MCNW.7}
\]
There is no additional local, torsion, critical or prime-count cost.
The same $O_\varepsilon(u^{806+16\varepsilon}(\log u)^2)$ envelope,
identical common window and preceding threshold/slack implications
remain valid.
\end{theorem}

\begin{proof}
\proofstep{1. The full split-cube theorem on the literal original coefficient.}
For an eligible prime $\ell$, the actual host contains the
Kummer cubic factor and all three quadratic characters in $\Delta_8$.
Their values require only $\mu_6$, contained in the prescribed
original completion $O_\ell$ of $\mathbb Z[\zeta_{3\cdot2^h}]$.
Keep its original fraction field, embedding, uniformizer and
normalizing rank $b_\ell$. The exact (FC.2) coefficient is
\[
 (O_\ell\otimes_{\mathbb Z_\ell}\mathbb Q_\ell/\mathbb Z_\ell)
                                      (\chi^j\delta)(-e)
       =(E_\ell/O_\ell)(\chi^j\delta c_\ell^{-e}).
 \tag{MCNW.8}
\]
Apply AC-G215 (MCN.2)--(MCN.4) with $\kappa=1$.
The prime in (MCNW.2) is split, $269$ is a cube modulo $\ell$,
and the selected quadratic sign is $-1$. These are exactly
the additional branch of its full split-prime criterion.
Its Step4 preserves original-host inflation and all full local
conditions. Thus the actual strict group in (MCNW.8) is zero.

The class input is the retained G207 certificate and targeted
class-sector argument, which applies before any local cube test.
Its certified smaller-field quotients have orders $7,63,42$;
their prime-to-$\ell$ conclusion is valid because $\ell>3$
and $\ell\ne7$. No new class computation or equality of actual
class numbers is used by this composition.

At each of the two coefficient places, $K_v=\mathbb Q_\ell$.
The cube condition and Hensel lifting split $X^3-269$ locally,
so the Kummer factor is trivial there. The selected quadratic
Frobenius has value $-1$, whose difference from one is an
original-$O_\ell$ unit. The full coefficient invariants vanish.
All away-inertia invariant quotients also vanish as in G215,
including the primes over $2,3,269$. Its original-uniformizer
argument therefore uses the same full conditions as (FC.2).
No rational noncritical replacement is needed for these mixed zeros.
The finite-etale dictionary on
$\operatorname{Spec}\mathcal O_K[1/(2\cdot3\cdot269\cdot\ell)]$
is retained from G215, without a new archival or motivic claim.

\proofstep{2. Exact old intersections and the complete prime-family count.}
Every new label has nontrivial Kummer factor.
All the rational families $V_u,Y_u,C_5^{\rm all}(u),C_7^{\rm all}(u)$
and every old rational family in $Z_u$ have trivial Kummer factor.
The old mixed $Z_u$ families have coefficient $269$, explicitly
excluded from $\mathscr C$.

The entire family $M_8(u)$ uses noncube primes, whereas $C(u)$
uses cube primes; thus their coefficient-prime sets are disjoint.
The family $M_7(u)$ may have a prime in common with $\mathscr C$,
but its rational-seven factor is a nontrivial cubic character.
The factor of each new label at seven is trivial.
Independence of the factors in (FC.1) separates them; prime
disjointness is neither needed nor claimed in this case.
Finally both $Q_7(u)$ and $Q_{7,7}(u)$ have coefficient seven,
whereas $\mathscr C$ excludes seven. In particular all the new
four-per-weight residual family from G214 stays disjoint.
This exhausts every old omission and proves (MCNW.4) and (MCNW.6).

For split primes different from three the residues modulo $24$
are $1,7,13,19$. The three quadratic signs in that order are
\[
 \begin{array}{c|rrr}
 \ell\bmod24&\chi_{-4}(\ell)&\chi_8(\ell)&\chi_{-8}(\ell)\\\hline
 1&1&1&1\\
 7&-1&1&-1\\
 13&1&-1&-1\\
 19&-1&-1&1
 \end{array}\tag{MCNW.9}
\]
Every prime of $\mathscr C$ therefore admits exactly two
quadratic labels and both Kummer powers, giving four labels.
There are $\lfloor N/(\ell-1)\rfloor$ positive eligible weights,
zero unless $\ell\leq T$. This proves $4J_C(T)$ with no
prime-distribution assumption. The preceding count (M77W.6)
then gives (MCNW.5)--(MCNW.6).
Since $269\equiv10\pmod{37}$, $10^{12}\equiv1\pmod{37}$ and
$37\equiv13\pmod{24}$, the four $\chi^j\chi_8,\chi^j\chi_{-8}$
labels at weight $36$ prove properness for $T\geq37$.
The entire prime set is used at once, rather than only this example.

\proofstep{3. Pointwise mixed deletion and the inherited cost.}
Each new summand is $\log1/b_\ell=0$ by Step1.
Remove those terms from (M77W.3), equivalently apply (MCN.8)
to the currently retained actual triples. This gives the exact
pointwise equality in (MCNW.4). The whole rational noncritical
sum, every old mixed omission and all critical partners are retained.
No positive derived index is split or asserted zero.

Substitute the equality into (M77W.7). One full (FL.4)
local majorant, one (LT.2) torsion aggregate, one existing
critical aggregate and the old G189 full-index cost still give
(MCNW.7). There is no multiplier for four, $J_C(T)$, a proof-field
degree, or the number of eligible primes.
The actual degrees $D_u=144\cdot2^h$ and $d_u=96\cdot2^h$
retain their distinct roles and the same error envelope.

\proofstep{4. Identical window, threshold and simultaneous quantifiers.}
Fix the same original datum, page, $0<\varepsilon<1/256$,
parameters $J,X_0,s$, prescribed pairs and common window
$X^{1/J}\leq W$, $W^2\leq X$ of (TC.3).
The two equal remainders have identical numerical premises,
thresholds and slack on this same $W$.
For the full original index retain the same $U_0$,
threshold $\max(X_0,U_0^J)$ and half-slack argument.
No independent existential windows are intersected.
All original fields, coefficient rings, lattices, finite conditions,
actual ranks $R.t$, pages and prescribed pairs remain under
the same simultaneous source quantifiers.
\end{proof}

\paragraph{Scope and chronology.}
The composition retains G214's full coefficient-seven residual family
and adds the entire split-cube quadratic complement in one step.
Every other unsupported mixed or noncritical term remains in (MCNW.3).
The remaining numerical bound, mixed theta/Selmer realization and
H04.4 remain open. G215's full arbitrary-set identity was already
assessed in Review32; this explicit sequential composition and its
later integration do not earn separate or banked movement credit.
No new question row or automatic gate closure is asserted.
